% 矮行星（综述）
% license CCBYSA3
% type Wiki

本文根据 CC-BY-SA 协议转载翻译自维基百科\href{https://en.wikipedia.org/wiki/Dwarf_planet}{相关文章}。

\begin{figure}[ht]
\centering
\includegraphics[width=6cm]{./figures/d21611a5210b90a2.png}
\caption{} \label{fig_Dwarf_1}
\end{figure}
“矮行星是一类直接绕太阳运行、具有行星质量但体积较小的天体。它们的质量足以使自身达到由重力塑形的近球形，但不足以像太阳系中八大经典行星那样在轨道上实现支配地位。矮行星的原型是冥王星，它曾被视为一颗行星长达数十年，直到 2006 年‘矮行星’概念被正式采用。许多行星地质学家认为矮行星与具有行星质量的卫星都应被视为行星[1]，但自 2006 年以来，国际天文学联合会（IAU）和许多天文学家已将它们排除在行星名录之外。

矮行星可能具有地质活动能力——这一预期在 2015 年“黎明号”对谷神星的探测和“新视野号”对冥王星的探测中得到证实。因此，它们成为行星地质学家特别关注的对象。

天文学家普遍同意至少有九个最大的候选体属于矮行星——按直径从大到小大致排列为：冥王星、厄里斯、妊神星、鸟神星、公公星、夸欧尔、塞德娜、谷神星和奥库斯。对于第十大候选体萨拉西娅仍存在相当的不确定性，因此可将其视为临界案例。在这十个天体中，已有两个被航天器访问过（冥王星和谷神星），另有七个至少拥有一颗已知卫星（厄里斯、妊神星、鸟神星、公公星、夸欧尔、奥库斯和萨拉西娅），这使得它们的质量得以确定，从而可估算其密度。质量和密度又可用于构建地球物理模型，以尝试判断这些世界的性质。只有塞德娜既未被探测器访问，也没有已知卫星，因此其质量难以准确估计。一些天文学家还将许多更小的天体也纳入矮行星的范畴[2]，但目前并无共识认为这些小天体很可能是矮行星。”**
\subsection{概念的历史}
\begin{figure}[ht]
\centering
\includegraphics[width=6cm]{./figures/0ce0512cdfe2e0c4.png}
\caption{} \label{fig_Dwarf_2}
\end{figure}